\section{Введение}
\label{sec:intro}
Мы рассматриваем минимизацию выпуклой функции на компактном выпуклом
множестве $Q$. Функция дважды непрерывно дифференцируема в окрестности
$Q$ и удовлетворяет условию гладкости
\[
 \|\nabla^2 f(x)\|\le L_0+L_1\|\nabla f(x)\|,\qquad x\in Q.
\]
Параметры $L_0,L_1$ и диаметр $Q$ входят в оценку числа запросов;
самому методу с подбором шага их значения не нужны.
Неизвестное минимальное значение не входит ни в условие гладкости,
ни в алгоритм. Точка $x^*$ обозначает минимум только в анализе.

Полный зазор Франка-Вульфа
$G(x)=\max_{s\in Q}\langle\nabla f(x),x-s\rangle$
в силу выпуклости удовлетворяет
$f(x)-f(x^*)\le G(x)$ \cite{jaggi}.
Поэтому $f(x)-G(x)$ даёт вычисляемую нижнюю границу минимума.
Оценка верна и для граничного решения с ненулевым градиентом.
Остановку определяет зазор, а не требование малой нормы градиента.

При этом зависимость кривизны от градиента влияет на выбор шага.
Мы показываем этот эффект на семействе выпуклых функций
\[
 f_A(x,y)=\exp(x)+Ay,\qquad
 Q=[-1,1]\times[0,1],\qquad A>0.
\]
Все функции имеют одни и те же параметры $L_0=L_1=1$.
При старте из $(0.7,0)$ полный линейный оракул всегда возвращает
$(-1,0)$. Зазор и значения функции вдоль посещаемой грани не зависят
от $A$, хотя норма градиента растёт вместе с $A$.
Мы доказываем, что шаг из алгоритма~1 Выгузова и Стонякина
\cite{vyguzov} с фиксированными параметрами гладкости требует
$\Omega(A)$ обновлений даже для пересечения прямой $x=0$.
При точных ответах оракула и одинаковой начальной оценке кривизны
квадратичный бэктрекинг даёт одни и те же пробные точки, принятые
шаги и число запросов до остановки при всех $A$.
Сравнение относится к их алгоритму~1 с фиксированными параметрами.
Оно не распространяется на адаптивный алгоритм~2.

Далее строится правило остановки с адаптивной абсолютной точностью
направленного оракула. Границы его ошибки подтверждают и полный
зазор, и величину спуска вдоль выбранного направления.
Запрошенная точность переносится между итерациями и уточняется при
необходимости. Алгоритм сохраняет лучшую подтверждённую нижнюю
границу минимума. Доказана остановка с заданной ошибкой по функции
и получены оценки всех вызовов оракула и вычислений функции.
Подсчёт включает неудачные проверки точности, отклонённые пробные
шаги и заключительную проверку остановки. Точные градиенты и точная
линейная минимизация по всему $Q$ входят в эту схему как частный случай.

\paragraph{Связь с известными результатами.}
Сертификаты через зазор и неточные линейные оракулы уже входят в
анализ выпуклого метода Франка-Вульфа \cite{jaggi}.
Выгузов и Стонякин \cite{vyguzov} изучают фиксированные и адаптивные
шаги при указанном условии гладкости.
Квадратичный бэктрекинг также известен \cite{pedregosa}.
Мы используем этот способ шага и доказываем приведённые выше гарантии
остановки с полным подсчётом запросов. Явное семейство количественно
показывает, как компонента градиента влияет на шаг через норму,
хотя выбранное направление и значения функции вдоль него не меняются.

Условие кривизны через невязку функции из анализа разогрева шага
\cite{warmup} является другой моделью и здесь не предполагается.
Метод Local LMO \cite{locallmo} минимизирует линейную функцию
по $Q\cap B(x,t)$, а наше правило сохраняет полный линейный оракул
по $Q$. Таким образом, используется исходная выпуклая постановка,
без регулярности относительно неизвестного центра и без знания
$f(x^*)$.
